%%%%%%%%%%%%%%%%%%%%%%%%%%%%%%%%%%%%%%%%%%%%%%%%%%%%%%%%%%%%%%%%%%%%%%%%%%%%%%%%

\pagestyle{empty}
\begin{abstract}
  Las interfaces cerebro-computadora (BCI) representan un paradigma prometedor para la interacción humano-máquina, pero su desarrollo enfrenta limitaciones significativas debido a la falta de herramientas accesibles que faciliten la captura sincronizada de señales electroencefalográficas (EEG) junto con eventos de entrada convencionales del usuario. Esta sincronización precisa es fundamental para establecer correlaciones entre actividad cerebral y acciones del usuario, requisito esencial para el entrenamiento de modelos predictivos en aplicaciones BCI.

  Este trabajo presenta el desarrollo de un \textit{framework} modular de código abierto que aborda esta problemática integrando la captura sincronizada de señales EEG y eventos de entrada de usuario (teclado, ratón, controladores) con precisión temporal de milisegundos. El sistema implementa una arquitectura basada en principios SOLID con patrones dirigidos por eventos, proporcionando múltiples interfaces de usuario (GUI PyQt5, CLI y API REST con WebSocket) que facilitan tanto el uso interactivo como la integración programática.

  La solución soporta dispositivos comerciales de EEG mediante BrainFlow, incluyendo el g.tec Unicorn Hybrid Black, e incorpora procesamiento de señales en tiempo real con algoritmos avanzados como análisis de componentes independientes (ICA), transformadas \textit{wavelet} y filtros digitales. El \textit{framework} incluye visualización de señales biomédicas, monitoreo de calidad de señal, y exportación a formatos estándar como BIDS, EDF+, HDF5 y CSV.

  La validación del sistema con datos sintéticos demuestra sincronización precisa entre señales EEG y eventos de entrada, cumpliendo los requisitos para la generación de \textit{datasets} de calidad orientados al desarrollo de interfaces cerebro-computadora. El \textit{framework} desarrollado constituye una herramienta completa que integra capacidades tradicionalmente fragmentadas en el ecosistema BCI, facilitando la investigación y desarrollo en este campo emergente.

  \vspace*{25pt}
  \begin{segundoresumo}
    Brain-computer interfaces (BCI) represent a promising paradigm for human-machine interaction, but their development faces significant limitations due to the lack of accessible tools that facilitate synchronized capture of electroencephalographic (EEG) signals along with conventional user input events. This precise synchronization is fundamental for establishing correlations between brain activity and user actions, an essential requirement for training predictive models in BCI applications.

    This work presents the development of an open-source modular \textit{framework} that addresses this problem by integrating synchronized capture of EEG signals and user input events (keyboard, mouse, controllers) with millisecond temporal precision. The system implements an architecture based on SOLID principles with event-driven patterns, providing multiple user interfaces (PyQt5 GUI, CLI, and REST API with WebSocket) that facilitate both interactive use and programmatic integration.

    The solution supports commercial EEG devices through BrainFlow, including the g.tec Unicorn Hybrid Black, and incorporates real-time signal processing with advanced algorithms such as Independent Component Analysis (ICA), \textit{wavelet} transforms, and digital filters. The \textit{framework} includes biomedical signal visualization, signal quality monitoring, and export to standard formats such as BIDS, EDF+, HDF5, and CSV.

    System validation with synthetic data demonstrates precise synchronization between EEG signals and input events, meeting requirements for generating quality \textit{datasets} oriented towards brain-computer interface development. The developed \textit{framework} constitutes a comprehensive tool that integrates traditionally fragmented capabilities in the BCI ecosystem, facilitating research and development in this emerging field.
  \end{segundoresumo}
\vspace*{25pt}
\begin{multicols}{2}
\begin{description}
\item [\palabraschaveprincipal:] \mbox{} \\[-20pt]
  \begin{itemize}
    \item Interfaces cerebro-computadora (BCI)
    \item Electroencefalografía (EEG)
    \item Sincronización de datos
    \item Framework modular
    \item Procesamiento de señales
    \item Captura de eventos
    \item Datasets biomédicos
    \item BrainFlow
    \item Arquitectura dirigida por eventos
    \item API REST
    \item PyQt5
    \item Python
  \end{itemize}
\end{description}
\begin{description}
\item [\palabraschavesecundaria:] \mbox{} \\[-20pt]
  \begin{itemize}
    \item Brain-computer interfaces (BCI)
    \item Electroencephalography (EEG)
    \item Data synchronization
    \item Modular framework
    \item Signal processing
    \item Event capture
    \item Biomedical datasets
    \item BrainFlow
    \item Event-driven architecture
    \item REST API
    \item PyQt5
    \item Python
  \end{itemize}
\end{description}
\end{multicols}

\end{abstract}
\pagestyle{fancy}

%%%%%%%%%%%%%%%%%%%%%%%%%%%%%%%%%%%%%%%%%%%%%%%%%%%%%%%%%%%%%%%%%%%%%%%%%%%%%%%%